\chapter{Introduction}
\pagenumbering{arabic} \setcounter{page}{1}

Ice, which is formed from the freezing of water, is one of the most ubiquitous substances on Earth, and yet it exhibits a remarkable diversity of molecular geometries depending on temperature and pressure\cite{Zheligovskaya2006}. In fact, scientists have identified at least 20 crystaline phases of ice, each with its own unique crystal structure and properties\cite{BartelsRausch2012}\cite{Fuentes-Landete}. There exists three amorphous phases of ice on top of the crystal ones\cite{Fuentes-Landete}. Understanding the characteristics of these different forms of ice is crucial for multiple scientific fields, including climate science and materials science, as it provides valuable insights into their behavior and properties.\\
The most familiar form of ice is Ice-Ih, which has a hexagonal crystal structure and interlocking hydrogen bonds between water molecules. This form of ice is what we typically encounter in our daily lives, such as in snowflakes and ice cubes\cite{Baker2019}. Conversely, Ice-II has a rhombohedral crystal structure and higher density than Ice-Ih and forms under high pressures and low temperatures\cite{Kamb1964}. Ice-III can exist at temperatures as low as $-100 \degree C$ and under greater pressures than Ice-II, making it useful in high-pressure experiments\cite{Bauer_2008}. Ice-V is one of the densest form of ice, characterized by its monoclinic crystal structure, forms at 253K temperature and 500MPa\cite{Hasted_1969}. Similarly, Ice-VI is denser than Ice-III and forms under high pressures and low temperatures, providing insights into planetary science\cite{Kuhs_1984}. Ice-VII is a cubic crystal form of ice with a unique crystal structure and properties that are still being studied\cite{Kuhs_1984}.\\
Cubic ice, or Ice Ic, assumes a cubic shape and has different hydrogen bonding than other forms of ice\cite{Murray_2006}. Clathrate structures, which form when gas molecules are trapped within the ice lattice, are crucial in the study of gas hydrates, the storage and release of methane from the ocean floor, and the Earth's climate\cite{Buffett2004}\cite{Buchanan2005}. The classification of these packing geometries has been achieved through considerable effort in molecular structure analysis. While conventional crystallographic analysis has been used, current molecular simulation shows promise in providing precise real-space geometrical information for many molecular systems. However, molecular geometry classification remains a challenging task. To efficiently distinguish between different structures, order parameters are often utilized\cite{CHAU_1998}\cite{Giovambattista_2005}\cite{Yan_2007}\cite{Zimmermann2017}\cite{Yan2007}\cite{Errington2001}\cite{Piaggi2017}\cite{Lechner2008}\cite{Reinhardt2012}\cite{Mickel2013}. Moreover, order parameters can be defined for local structures to describe transient structures during phase transitions.\\
The structure of ice is of great interest as it exhibits a wide range of fascinating properties, such as high thermal stability, unique optical and electronic properties, and important implications in climate science. The regular lattice of water molecules forming a tetrahedral arrangement through hydrogen-bond interactions characterizes the structure of ice, while the local structure of liquid water lacks long-range order, and only a partial short-range order remains, influenced by various factors such as pressure, temperature, solutes, and hydrophobic groups\cite{Galamba_2013}. By gaining a deeper understanding of the molecular geometries of ice, scientists can better understand the behavior and properties of this vital substance, with potential applications in fields such as materials science, energy storage, and climate modeling.\\
The intricate phase space of ice and water gives rise to a multitude of transitions across various phases\cite{Hansen2021}\cite{Chaplin2019}\cite{Bore2022}\cite{Hirata2017}. These transitions encompass a wide spectrum of direct-indirect and ordered-disordered phenomena, which poses a considerable challenge in molecular simulation due to the vastness of the configuration space. The need for a mathematical functional, specifically an order parameter, becomes imperative to effectively describe these phenomena and capture the nuanced microscopic events occurring.\\
However, this task becomes considerably complex owing to the subtle differentiations existing across different phases. Sometimes, the similarities are so intricate that they cannot be discerned using simple geometric functions\cite{Moore2011}. While researchers have leveraged their domain knowledge to devise such order parameters, they tend to be highly specific to individual cases and often require computationally intensive calculations.

\begin{figure}[H]
\centering
\includegraphics*[width=.85\linewidth]{intro/phase.png}
  \caption{Phases labeled with black Roman numerals are hydrogen-disordered phases; those labeled with red numerals are hydrogen-ordered phases; those labeled with blue numerals are metastable phases. The red solid arrows represent direct disorder-order transitions. The red dotted arrows represent disorder-order transitions that cross the stability range of another phase\cite{Hansen2021}.}
\label{fig:Phase diagram}
\end{figure}
Among the widely utilized order parameters is the tetrahedral order parameter, extensively employed to probe local structures around water molecules and their hydrogen bonding networks. Ranging from 1 for a pure tetrahedral network to approaching 0 for random configurations, this parameter offers insight into the local environment\cite{CHAU_1998}\cite{Galamba_2013}\cite{Giovambattista_2005}\cite{Yan2007}.
In distinguishing between clathrates and bulk water, the four-body structural order parameters ($F_4$) have gained popularity. By considering the dihedral formed by two neighboring oxygen atoms (within a defined cutoff) and two outward hydrogen atoms attached to the oxygen molecules,the $F_4$ parameters provide valuable discrimination\cite{Rodger1996}.
Recent trends highlight the utility of entropy fingerprints and their localized counterpart, the local entropy\cite{Piaggi2017} . By employing the radial distribution function, these methods capture the fluctuation of conformations with significant success.
An intuitive approach to address this challenge traditionally revolves around the introduction of bond orientational order parameters. These parameters, based on rotationally invariant spherical harmonics, have found success in describing crystallization events. Steinhardt's order parameters, in particular, have been instrumental in distinguishing various ice phases based on their characteristic order parameter distributions\cite{Plimpton1995}.
A more intricate approach involves the utilization of atom-centered symmetry functions\cite{Behler2011}. With a unique atomic descriptor generated for the distribution of particles in 3D, the Smooth Overlap of Atomic Positions (SOAP) technique stands out\cite{De2016}\cite{Maksimov2020}. SOAP offers a local descriptor that meticulously encapsulates the subtle details of the immediate spatial environment surrounding a specific point. Through Gaussian smearing and transformation into coefficients using an orthonormal basis function system, SOAP effectively mitigates redundancies associated with rotation, translational, and permutation\cite{Himanen2020}.\\
However, despite these efforts, the design of physical intuition-based order parameters remains challenging. They tend to be specific to particular systems, computationally demanding, and may sometimes overlook valuable concealed information\cite{Donkor2023}\cite{Fulford2019}\cite{Doi2021}. In response to these challenges, the last decade has witnessed a surge in the utilization of Machine Learning (ML) methods. Data-driven models are now being employed to identify various local patterns present in simulation snapshots. Leveraging the capacity of deep neural networks to uncover hidden patterns within seemingly random distributions has paved the way for a new era in order parameter exploration..   
In recent years, the landscape of order parameter exploration has been dramatically expanded by the advent of various Machine Learning (ML) approaches\cite{ML_Method1}\cite{Cheng2019}. Convolutional Neural Networks (CNNs) have demonstrated remarkable proficiency in capturing spatial patterns within simulation snapshots\cite{Li2022}. These networks, often utilized in image recognition tasks, have been adapted to analyze the structural arrangements of molecules.
Graph Neural Networks (GNNs) provide an ideal framework for dealing with molecular systems, where atoms and their interactions can be represented as nodes and edges within a graph. GNNs excel at capturing complex dependencies and interactions, making them particularly suited for identifying structural motifs within simulations\cite{Hofstetter2022}\cite{Kim2020}.
Autoencoders, a type of neural network architecture, facilitate the extraction of lower-dimensional representations from high-dimensional data. They are adept at discerning latent features and structural characteristics that may not be immediately apparent in the original data. By training on simulation snapshots, autoencoders can uncover hidden relationships and patterns within the molecular configurations\cite{Kim2020}.
Incorporating these ML methods has revolutionized the study of order parameters, enabling the discovery of intricate correlations and complex structural changes that were previously elusive. As our understanding of ML techniques deepens, they promise to further enrich our insights into the phase transitions and behaviors of ice and water, unveiling hidden intricacies that were once beyond our grasp.\\
In this study, the local structure of various ice phases is investigated, and several order parameters are calculated to characterize them. A range of geometric order parameters, including orientational and translational tetrahedral order parameters, local structure index, local density based on Voronoi polyhedra, asphericity of the Voronoi cell, water-water angular distribution function, and Steinhardt’s bond orientational order parameter etc. have been proposed to quantify the local structure of ice phases. Our aim is to provide a comprehensive understanding of the local structure of different ice phases and the performance of various order parameters in characterizing them.


